%% CVEX v3 technical report -- ACM CCS (acmart sigconf) style.
%% Build: python3 gen_stats.py && make
\documentclass[sigconf,nonacm]{acmart}

\settopmatter{printacmref=false}
\setcopyright{none}
\renewcommand\footnotetextcopyrightpermission[1]{}
\pagestyle{plain}

\usepackage{xspace}
\usepackage{listings}
\usepackage{booktabs}
\usepackage{multirow}
\usepackage{tikz}
\usetikzlibrary{positioning,fit,arrows.meta,backgrounds}

% Generated by gen_stats.py from CVEX_v3-records -- do not edit by hand.
\newcommand{\NumRecords}{24\xspace}
\newcommand{\NumFunctional}{10\xspace}
\newcommand{\NumScaffolds}{14\xspace}
\newcommand{\NumQuarantined}{0\xspace}
\newcommand{\NumGenieScaffolds}{0\xspace}
\newcommand{\NumGenieFunctional}{0\xspace}
\newcommand{\NumMigratedFunctional}{10\xspace}
\newcommand{\NumFlag}{9\xspace}
\newcommand{\NumVerifier}{0\xspace}
\newcommand{\NumCrash}{1\xspace}
\newcommand{\NumChanStdout}{5\xspace}
\newcommand{\NumChanNetwork}{4\xspace}
\newcommand{\NumGenieCWEs}{0\xspace}
\newcommand{\NumGenieProjects}{0\xspace}
\newcommand{\PcapMiB}{0.1\xspace}
\newcommand{\SyscallMiB}{33\xspace}


\newcommand{\cvex}{CVEX\xspace}
\newcommand{\cvexthree}{CVEX~v3\xspace}
\newcommand{\genie}{CVE-Genie\xspace}
\newcommand{\file}[1]{\texttt{#1}}

\lstdefinestyle{cvex}{
  basicstyle=\ttfamily\scriptsize,
  columns=fullflexible,
  keepspaces=true,
  breaklines=true,
  frame=single,
  framerule=0.4pt,
  rulecolor=\color{black!40},
  backgroundcolor=\color{black!3},
  commentstyle=\color{black!55}\itshape,
  showstringspaces=false,
  xleftmargin=2pt,
  xrightmargin=2pt,
  aboveskip=6pt,
  belowskip=4pt,
}
\lstset{style=cvex}

\begin{document}

\title{CVEX v3: Executable CVEs for Verifiable Vulnerability Reproduction}
\subtitle{Technical Report}

\author{Giovanni Vigna}
\affiliation{%
  \institution{University of California, Santa Barbara}
  \city{Santa Barbara}
  \state{CA}
  \country{USA}}
\email{vigna@ucsb.edu}

\begin{abstract}
A CVE entry tells a defender that a vulnerability exists, but not how to
make it happen. Reproducing a vulnerability, that is, rebuilding the
vulnerable software, configuring it, and running an exploit that provably
succeeds, remains a slow, expert-driven process, and the result is rarely
packaged in a form that others can rerun. This is a problem for anyone who
needs ground truth about attacks: the developers of detection models, the
analysts who assess whether a patch actually fixes a flaw, and the
researchers who build and evaluate automated security tools. We present
\cvexthree (\emph{executable CVE}), a standard format for describing the
reproduction of a vulnerability as a self-contained, containerized record.
A \cvex record bundles the vulnerable target, a benign workload, an
exploit, a recorder, and a machine-checkable success criterion. Borrowing
from Capture-the-Flag competitions and from the verifiers produced by the
\genie reproduction framework, success is defined by the capture of a
secret flag that only a working exploit can reach. Every record is run
twice, once with benign traffic and once with the attack, and both runs are
captured as packet traces and system-call traces, which yields a labeled,
paired dataset. We describe the format and its runner, and release a
corpus of \NumRecords records, \NumFunctional of which are fully functional and
verified, with reference datasets.
\end{abstract}

\keywords{vulnerability reproduction, CVE, exploits, datasets, containers,
intrusion detection, benchmarks}

\maketitle

\section{Introduction}
\label{sec:intro}

\subsection{The Problem of Vulnerability Reproduction}

The Common Vulnerabilities and Exposures (CVE) program assigns an identifier
and a short description to every publicly disclosed vulnerability, and the
number of entries grows by tens of thousands every year~\cite{cve,nvd}. A CVE
entry is a pointer, not an experiment: it names the affected product and
versions, and it may link to an advisory or a patch, but it does not tell the
reader how to bring the vulnerable software up, how to trigger the flaw, or
how to tell whether the attempt worked. Turning an entry into a working
reproduction is left to whoever needs it.

Many people need it. Reproductions are the ground truth on which much of
security research and practice rests:

\begin{itemize}
\item \textbf{Detection models.} Network and host intrusion detection systems,
signature generators, and machine-learning anomaly detectors must be trained
and evaluated on traffic and system activity that contains real, successful
attacks, alongside comparable benign activity. Long-lived public datasets are
either synthetic or stale, and evaluating detectors on them has been
repeatedly criticized~\cite{sommer2010outside,lippmann2000darpa,creech2013adfa}.
A reproduction that can be rerun at will produces fresh, labeled captures for
exactly the vulnerability of interest.

\item \textbf{Patch analysis.} Security patches are often incomplete:
large-scale studies show that a notable fraction of fixes must be revised
before a vulnerability is really closed~\cite{li2017patches}. The two Apache
\texttt{httpd} path-traversal vulnerabilities in our corpus illustrate this:
the fix for CVE-2021-41773 shipped in version 2.4.50 was itself bypassable,
which led to CVE-2021-42013. An executable reproduction turns ``does the
patch work?'' into a test that can be run against the patched build, and
against every later version.

\item \textbf{Automated security tools.} Fuzzers, program analyzers, exploit
generators, and, more recently, LLM-based agents are all evaluated on
vulnerabilities whose presence and exploitability must be established
beforehand~\cite{hazimeh2020magma,zhu2025cvebench,wang2026cybergym}. Without
reproducible environments, these evaluations are hard to compare and hard to
repeat.

\item \textbf{Risk assessment and education.} Whether a vulnerability is
exploitable in a given configuration, and what an exploitation looks like on
the wire and on the host, are questions that defenders answer most reliably
by running the attack. Reproducing attacks is also one of the most effective
ways to teach how they work.
\end{itemize}

Reproduction, however, is expensive. In the most extensive study of the
subject, a team of analysts spent 3,600 hours trying to reproduce 368
crowd-reported memory-corruption vulnerabilities, and found that reports
routinely omit the information needed to succeed~\cite{mu2018reproducibility}.

\subsection{The Challenge: No Standard Reproduction Environment}

The difficulty is not only in crafting an exploit. Most of the effort goes
into the environment: obtaining the exact vulnerable version, building it
with a compatible toolchain, installing its dependencies, configuring it into
an exploitable state, and seeding it with the data an attack needs. Every
reproduction effort solves these problems in its own way, and the solution is
rarely captured in a reusable form. Exploit databases publish exploits but
not the environment they target~\cite{exploitdb}. Collections of
containerized vulnerable applications, such as Vulhub~\cite{vulhub}, provide
environments but no uniform notion of what a successful exploitation is, and
no instrumentation. Benchmarks built for fuzzers or AI agents define
environments tailored to their own harnesses~\cite{mei2026arvo,zhu2025cvebench,wang2026cybergym,lee2025secbench}.
As a result, there is no standard way to \emph{describe} the environment in
which a vulnerability is reproduced, to \emph{verify} that the reproduction
succeeded, or to \emph{record} what happened during the attack. Results
produced by one group cannot easily be rerun by another, and reproductions
are lost when the original setup disappears.

\subsection{Our Solution: CVEX, the Executable CVE}

We propose \cvex, short for \emph{executable CVE}: a standard format that
augments a CVE entry with everything needed to reproduce it. A \cvex record
is a small, self-contained directory that describes (i)~the vulnerable
\emph{target}, built in a container; (ii)~a \emph{benign} workload that uses
the same service legitimately; (iii)~an \emph{attack} that exploits the
vulnerability; (iv)~a \emph{recorder} that captures network packets and
system calls; and (v)~a machine-checkable \emph{success criterion}. Running a
record produces a verdict and a dataset.

The central design choice of the third version of the format, \cvexthree, is
how success is defined. Following Capture-the-Flag (CTF) competitions, and
the verifiers produced by the \genie automated reproduction
framework~\cite{ullah2026cvegenie}, a record plants a secret \emph{flag} in
the target, at a location that only a successful exploitation can reach: a
root-owned file for a privilege escalation, a file outside the web root for a
path traversal, a database row for an SQL injection. The attack succeeds if
and only if it returns the flag. This makes verification objective and
uniform across vulnerability classes, and it guarantees that the attack half
of every dataset contains an exploitation that actually worked.

\subsection{Contributions}

This report makes two contributions.

\begin{enumerate}
\item \textbf{The \cvexthree format.} We define a uniform, Docker-based record
structure with a declarative manifest, three deployment topologies, and three
success criteria (planted flag, CTF-style verifier, and crash witness). We
define a run protocol that exercises every record in two isolated phases,
benign and attack, each against a pristine target, and records both as packet
captures and system-call traces. We provide a reference runner that builds a
record, runs it, verifies it, and writes a content-addressed dataset
(Section~\ref{sec:format}).

\item \textbf{A corpus of executable CVEs.} We release \NumRecords \cvexthree
records. \NumFunctional of them are fully functional: they have been built
and run by the reference runner, their exploitation has been verified, and
they ship with reference benign/attack captures. These include \NumMigratedFunctional
records migrated by hand from the previous generation of the format
and \NumGenieFunctional records converted from \genie reproductions, covering
\NumGenieCWEs CWE classes across \NumGenieProjects open-source projects. The
remaining records are scaffolds that are being functionalized
(Section~\ref{sec:dataset}).
\end{enumerate}

The rest of this report is organized as follows. Section~\ref{sec:background}
reviews CVEs, CWEs, existing approaches to reproduction, and the earlier
versions of \cvex. Section~\ref{sec:format} describes the \cvexthree format,
its runner, and the corpus. Section~\ref{sec:examples} walks through two
records in detail. Section~\ref{sec:future} discusses extending the format to
virtual machines, Section~\ref{sec:related} surveys related work, and
Section~\ref{sec:conclusion} concludes.

\section{Background}
\label{sec:background}

\subsection{CVE and CWE}

The CVE program, operated by MITRE on behalf of the U.S.\ government, assigns
a unique identifier of the form \texttt{CVE-YYYY-NNNNN} to each publicly
disclosed vulnerability, together with a short description and a list of
references~\cite{cve}. The National Vulnerability Database (NVD) enriches CVE
entries with severity scores, affected-product configurations, and a weakness
classification~\cite{nvd}. The classification uses the Common Weakness
Enumeration (CWE), a hierarchical catalog of software and hardware weakness
types, such as CWE-22 (path traversal), CWE-79 (cross-site scripting), or
CWE-269 (improper privilege management)~\cite{cwe}. CVE and CWE answer the
questions \emph{what} is vulnerable and \emph{what kind} of flaw it is; they
deliberately do not describe \emph{how} to exploit it. A typical entry
consists of a few sentences of prose and links to an advisory, a patch
commit, and occasionally a proof of concept (PoC).

\subsection{Approaches to Vulnerability Reproduction}

\paragraph{Manual reproduction.}
The traditional approach is manual. An analyst reads the advisory and the
patch, locates and builds the vulnerable version, configures it, and writes
or adapts an exploit, iterating until the exploit works. Mu et
al.~\cite{mu2018reproducibility} quantified the cost of this process: a group
of security analysts spent 3,600 hours attempting to reproduce 368
vulnerabilities, and the information in individual reports was often
insufficient for success. Their survey of 43 practitioners found that
reproduction relies heavily on manual debugging and guesswork. Exploit
databases~\cite{exploitdb} and environment collections such as
Vulhub~\cite{vulhub} reduce part of this effort, but each covers only half of
the problem: one provides exploits without environments, the other
environments without a uniform notion of success or instrumentation.

\paragraph{Automated reproduction: CVE-Genie.}
\genie~\cite{ullah2026cvegenie}, to appear at ACM CCS 2026, automates the
whole process with a multi-agent system built on large language models.
Starting from a CVE identifier, \genie collects the associated resources
(source code at the vulnerable version, advisory, CWE, patch, and PoC when
available), rebuilds the vulnerable project, develops an exploit, and
produces a \emph{verifier}: a script that runs the exploit and checks its
effect. The system is organized in four stages, each made of
developer agents paired with critic agents: a Processor, a Builder, an
Exploiter, and a CTF Verifier. The verifier is designed after Capture-the-Flag
challenges: it prints a fixed flag string only if it observes that the
exploit had its intended effect, so that reproduction success can be checked
mechanically. Evaluated on 841 CVEs published between June 2024 and May 2025,
\genie reproduced 428 of them (about 51\%) across 267 projects, 141 CWEs, and
22 programming languages, at an average cost of \$2.77 per CVE. The authors
release the reproduced CVEs together with the full logs of the agents'
interactions.

\genie shows that reproduction can be automated at scale, but its outputs are
tied to the environment in which the agents worked: a long-lived machine on
which the agents installed packages, compiled code, wrote helper files, and
started services interactively. The reproduced artifacts (the source
snapshot, the exploit, and the verifier) do not capture this state, so they
cannot be rerun elsewhere without reconstructing it. Packaging reproductions
so that they can be rerun anywhere is precisely the problem that \cvex
addresses, and \genie's flag-based verifier is one of the main inspirations
for the success criterion of \cvexthree.

\subsection{Previous Versions of CVEX}

\cvexthree is the fourth iteration of an effort that began at UC Santa
Barbara in 2021. Table~\ref{tab:generations} summarizes the earlier versions.

\begin{table*}[t]
\centering
\caption{Evolution of the CVEX format.}
\label{tab:generations}
\small
\begin{tabular}{@{}lllll@{}}
\toprule
 & \textbf{v1 (2021--22)} & \textbf{v2 (2023--24)} & \textbf{VM framework (2024--25)} & \textbf{v3 (2026)} \\
\midrule
Isolation & Docker Compose & Docker Compose & Vagrant VMs + Ansible & Docker (+ optional services) \\
Roles & target, exploiter, collector & server, client, listener & target VM(s), attacker VM, router VM & target, benign, attacker, recorder \\
Descriptor & \texttt{docker-compose.yml} + \texttt{cvex.json} & \texttt{docker-compose.yml} & \texttt{cvex.yml} + playbooks & \texttt{cvex.yml} + Dockerfiles \\
Success check & file created/read, leak, \texttt{dmesg} & manual inspection & manual inspection & planted flag, verifier, crash witness \\
Observability & pcap & pcap & pcap, TLS (mitmproxy), strace, Procmon & pcap + strace, benign \emph{and} attack \\
Platforms & Linux & Linux & Linux and Windows & Linux \\
Entries & 52 & 6 & 24 & \NumRecords (\NumFunctional functional) \\
\bottomrule
\end{tabular}
\end{table*}

\paragraph{Version 1 (2021--2022).}
The first version defined a \cvex entry as a Docker Compose project with
three services: a \emph{target} running the vulnerable software, an
\emph{exploiter} running the attack, and an optional \emph{collector}, a
\texttt{tcpdump} sidecar sharing the target's network namespace. A JSON
metadata file, validated by a schema, listed the CVEs, the authors, and an
\texttt{exploit\_verification} clause stating how to decide success: a file
created or read in the target, a leaked address, or a kernel log message.
About fifty entries were produced, and a community scoreboard rewarded
contributors with points for CVEs that others had upvoted. Images were
served from a project registry, which has since become unavailable, so most
of these entries can no longer be run as published.

\paragraph{Version 2 (2023--2024).}
A team of undergraduate researchers revisited the format with a simpler
server/client/listener Compose pattern, published prebuilt images on the
GitHub container registry, and catalogued the entries on a website. The
version dropped the metadata file and the formal success check: success was
established by inspecting the captured traffic.

\paragraph{The VM-based framework (2024--2025).}
The following version replaced containers with full virtual machines managed
by Vagrant and provisioned by Ansible, in order to support OS-level and
Windows vulnerabilities. Reusable \emph{blueprints} describe a network of
VMs (e.g., an Ubuntu attacker and an Ubuntu target), and per-CVE
\emph{records} add the playbooks and the attack commands. A router VM
captures all traffic, decrypts TLS with a transparent \texttt{mitmproxy},
and the framework traces system calls with \texttt{strace} on Linux and API
calls with Process Monitor on Windows. Twenty-four records were contributed
through pull requests. The framework is powerful, but heavy: every record
requires multi-gigabyte VM images and a hypervisor, and, like its
predecessors, it has no formal definition of success. Whether an exploit
worked is left to the person inspecting the traces.

\paragraph{Lessons.}
Three lessons from these versions shaped \cvexthree. First, an entry must be
rebuildable from what it contains: entries that depended on a registry or on
prebuilt images disappeared with them. Second, success must be defined by the
record itself and checked by a machine; otherwise a dataset may contain
failed attacks labeled as attacks. Third, the format must be light enough to
run hundreds of records on a single host, which favors containers over VMs
for the large class of vulnerabilities that do not need kernel-level control.

\section{CVEX v3}
\label{sec:format}

\cvexthree has three defining properties. First, it is \emph{Docker-only}: a
record is a Docker Compose project, with no virtual machines, hypervisor, or
per-host provisioning. Second, every record is exercised twice, once with
legitimate activity and once with the attack, and both runs are captured as
network packets and system calls. Third, success is \emph{flag-based}: the
record plants a secret that only a successful exploitation can reach, and
the attack succeeds if and only if it retrieves it. The format descends from
the previous generations described in Section~\ref{sec:background}: it keeps
the capture sidecar and the explicit exploit verification of the first
generation, the server/client container pattern of the second, and the
\file{cvex.yml} manifest and process tracing of the VM-based framework. The
flag and the benign/attack pairing are new.

\subsection{Record Structure}

Every record has the same shape, so that tools and researchers can treat all
records uniformly. A record is a directory named after its CVE identifier:

\begin{lstlisting}
CVE-YYYY-NNNNN/
  cvex.yml            # manifest (Sec. 3.3)
  compose.yml         # derived Docker Compose project
  target/
    Dockerfile        # builds the vulnerable software
    vendor/           # pinned sources and binaries
    plant-flag.sh     # installs the flag (Sec. 3.4)
  benign/
    Dockerfile
    interact.sh       # legitimate use of the service
  attacker/
    Dockerfile
    exploit.sh        # the attack; must emit the flag
  recorder/           # shared capture image
  _source/            # provenance: original artifacts
  dataset/            # committed reference run (Sec. 3.7)
\end{lstlisting}

A record is composed of four container roles, each with a single
responsibility (Table~\ref{tab:roles}). A record may add a second target
container, typically a database, when the service needs one; everything else
stays fixed. The \file{\_source} directory preserves the artifacts from which
the record was derived (for example, the previous-generation record or the
\genie reproduction, together with the CVE metadata), so that the provenance
of every record is explicit.

\begin{table}[t]
\centering
\caption{Container roles in a \cvexthree record.}
\label{tab:roles}
\small
\begin{tabular}{@{}lp{0.68\columnwidth}@{}}
\toprule
\textbf{Role} & \textbf{Responsibility} \\
\midrule
\texttt{target} & runs the vulnerable software and holds the planted flag \\
\texttt{benign} & a client that uses the service as a legitimate user would \\
\texttt{attacker} & a client that exploits the service and retrieves the flag \\
\texttt{recorder} & captures packets and system calls during each phase \\
\bottomrule
\end{tabular}
\end{table}

The roles are arranged in one of three \emph{topologies}. In the
\texttt{target-attacker} topology, the benign and attacker clients are
separate containers that reach the target over a private network; in
\texttt{target+db-attacker}, the target is backed by a database container.
In the \texttt{single} topology, used for local vulnerabilities and for
libraries and command-line tools that have no network service, the benign
workload and the attack run inside the target container itself, as
\file{target/benign.sh} and \file{target/attack.sh}.

\subsection{Design Rationale}

Two choices shape the rest of the format. The first is that a single, uniform
notion of success is more valuable than a per-record test. A record whose
exploit ``ran'' tells a user very little: the exploit may have failed
silently, or the environment may not have been vulnerable at all. By
requiring every attack to produce a value that only a successful
exploitation can obtain, the format guarantees that the attack half of every
dataset is a genuine, successful exploitation, and makes verification a
simple string comparison that is identical across vulnerability classes.

The second choice is to separate the benign and attack phases rather than
mixing them. Running the two phases against identical, pristine targets
makes the two captures directly comparable, so that the difference between
them is the attack itself, and gives clean labels. A \emph{mixed} phase, in
which benign activity runs concurrently with the attack, can be added later
without changing the schema.

\subsection{The Manifest}

The \file{cvex.yml} manifest declares everything the runner needs: the
topology, the network layout, the flag and the success criterion, the
container images, the processes to trace, and the two phases.
Listing~\ref{lst:manifest} shows a manifest for a demonstration record that
serves a benign \texttt{nginx} instance.

\begin{lstlisting}[float=t,caption={A \file{cvex.yml} manifest (demonstration record).},label={lst:manifest}]
cvex_version: 3
cve: CVE-0000-00001
topology: target-attacker # | single | target+db-attacker
network:
  subnet: 10.13.0.0/24    # fixed addresses in the pcap
flag:
  value: "CVEX{0000000000000000000000000000demo}"
  plant: file:/opt/secret/flag
success: flag             # flag | verifier | crash
capture:
  flag_channel: network   # stdout | network | verifier
  tls_intercept: true
  tls_ports: [443]
services:
  target:
    image: cvex/cve-0000-00001-target@sha256:<digest>
    ip: 10.13.0.10
    trace: nginx          # process(es) to trace (regex)
    ready: "tcp:443"      # readiness condition
phases:
  benign:
    client: benign        # runs benign/interact.sh
    image: cvex/cve-0000-00001-benign@sha256:<digest>
    ip: 10.13.0.30
  attack:
    client: attacker      # runs attacker/exploit.sh
    image: cvex/cve-0000-00001-attacker@sha256:<digest>
    ip: 10.13.0.20
limits:
  wall_clock_seconds: 600
\end{lstlisting}

The \texttt{ready} condition tells the runner when the target is ready for a
phase to start: either a TCP port that accepts connections or a file that the
target creates when its initialization completes. The \texttt{trace} field is
a regular expression that selects the processes whose system calls are
recorded. Records converted from \genie additionally carry an
\texttt{origin} field and the upstream \texttt{source} and
\texttt{sw\_version} of the vulnerable software.

The format imposes a small set of rules, checked at review time: every image
is pinned by an \texttt{@sha256} digest; either \texttt{flag.value} and
\texttt{flag.plant} are present, or the record declares a different success
criterion; both phases are present; and addresses are static, so that
captures are reproducible.

\subsection{Success Criteria}

\paragraph{The flag.}
A flag is a short token of the form \texttt{CVEX\{}\textit{32 hex
digits}\texttt{\}}, planted inside the target at a location, or behind a
protection, that only the vulnerability can bypass. The secrecy of the flag
comes from \emph{access control}, not from its being unknown: a legitimate
user of the service cannot reach it, the benign client never sees it, and
only a successful exploitation returns it. Because the flag is gated by
access rather than by randomness, it can be a fixed value committed with the
record, which keeps the reference dataset byte-reproducible. A runner may
optionally generate a fresh flag for each run, to guard against an attack
that simply hard-codes the expected value; the gating is what proves the
exploitation either way.

The location of the flag follows the class of the vulnerability, using a
small set of standard placements (Table~\ref{tab:flags}). The
\texttt{plant} field of the manifest names the location and, where relevant,
the file mode, and \file{plant-flag.sh} installs the flag at container start.

\begin{table}[t]
\centering
\caption{Standard flag placements by vulnerability class.}
\label{tab:flags}
\small
\begin{tabular}{@{}p{0.36\columnwidth}p{0.56\columnwidth}@{}}
\toprule
\textbf{Class} & \textbf{Flag location} \\
\midrule
Command execution / injection & \file{/flag}, readable only by running a command on the target \\
Path traversal / file read & a file outside the served directory \\
File upload leading to execution & \file{/flag}, readable once uploaded content is executed \\
SQL injection & a table row that the application never exposes \\
Privilege escalation & \file{/root/flag}, owned by root, mode 0600 \\
Local code execution / deserialization & \file{/flag} or an environment value not exposed by the normal interface \\
Memory corruption (crash only) & none; a crash witness is used instead \\
\bottomrule
\end{tabular}
\end{table}

\paragraph{Capture channels.}
The flag returned by the attack is detected on the channel that the record
declares. With the \texttt{stdout} channel, the attack prints the flag, and
the runner extracts it from the attack's output; this is the universal
fallback, and the natural choice for local records. With the
\texttt{network} channel, the runner searches for the flag in the attack
packet capture, which also proves that the flag crossed the wire.

\paragraph{The verifier.}
Records converted from \genie use the \texttt{verifier} criterion. \genie
produces, for each reproduced CVE, a verifier program in the style of a CTF
checker: it sets up the vulnerable condition, runs the exploit, checks for
the effect of the vulnerability, and prints a fixed flag if and only if the
check passes. A \cvexthree record that declares \texttt{success: verifier}
runs this program inside the target during the attack phase, and the runner
accepts the reproduction if the verifier exits successfully and prints the
flag. The verifier criterion extends the flag idea to vulnerabilities whose
effect is not the disclosure of a secret, for example a parser that splits
one request into two.

\paragraph{The crash witness.}
Some vulnerabilities, such as memory-corruption bugs that are reproduced only
up to a crash, have no secret to read. For these, the record declares
\texttt{success: crash}, and the flag is replaced by a \emph{crash witness}:
the recorder confirms that a traced process terminated because of a signal,
or produced a core dump, during the attack phase and not during the benign
phase. This is the only deviation from the flag rule, and it is declared
explicitly.

In the current corpus, \NumFlag functional records use a planted flag
(\NumChanStdout on the \texttt{stdout} channel and \NumChanNetwork on the
\texttt{network} channel), \NumVerifier use a \genie verifier, and
\NumCrash uses a crash witness.

\subsection{Run Protocol}

A record is run by the reference runner in two isolated phases, each starting
from a \emph{pristine} target: the target containers are created fresh for
each phase, so that neither run can contaminate the other.

\begin{enumerate}
\item \textbf{Build.} The runner builds (or pulls) the images.
\item \textbf{Benign phase.} The runner starts the target (and database, if
any) and the recorder, plants the flag, waits for the \texttt{ready}
condition, starts the capture, runs the benign workload, stops the capture,
and tears the containers down.
\item \textbf{Attack phase.} The runner starts a fresh target and recorder,
plants the flag, waits for readiness, starts the capture, runs the attack,
collects the flag it emits, stops the capture, and tears the containers
down.
\item \textbf{Verify.} The runner compares the observed flag with the planted
one (or evaluates the verifier or the crash witness) and writes the dataset
and the verdict.
\end{enumerate}

The benign workload performs legitimate interactions with the same service
and, where possible, with the same endpoints that the attack uses: a normal
login, a normal search, a normal file request, a normal invocation of a tool.
The two captures are therefore directly comparable: same images, same
addresses, same target state at the start of the phase. Each phase is bounded
by the wall-clock limit declared in the manifest.

A second command reruns a record and compares the fresh dataset with the
committed one, checking that the flags match and comparing the traffic after
normalizing timestamps and ephemeral ports. Once a record has been verified in
this way, it is locked: a lock file records the hash of the verified record,
so that later modifications are detected.

\subsection{The Recorder}

The recorder is a single shared image. In \emph{sidecar} mode, it shares the
network namespace and the process namespace of the target container. It runs
\texttt{tcpdump} on the target's interfaces, including the loopback
interface, so that it also captures traffic between processes inside the
target, and it attaches \texttt{strace} to the processes selected by the
\texttt{trace} expression, following their children. In \emph{gateway} mode,
intended for multi-host records that use TLS, the recorder acts as the
default gateway of the record's network, captures traffic on the bridge, and
runs a transparent \texttt{mitmproxy} whose certificate authority is
installed in the client and target images at build time, so that encrypted
traffic can be decrypted reproducibly.

\subsection{Dataset Layout}
\label{sec:dataset}

The output of a run is a self-contained dataset directory:

\begin{lstlisting}
dataset/
  benign/
    benign.pcap          # legitimate traffic
    benign.syscalls.log  # system calls under normal use
  attack/
    attack.pcap          # attack traffic
    attack.syscalls.log  # system calls during the attack
  flag.observed          # the flag the attack retrieved
  verdict.json           # criterion, channel, planted vs.
                         # observed flag, success
  MANIFEST.sha256        # content hashes of all files
\end{lstlisting}

For local records the packet captures may be empty or contain only loopback
traffic, and the system-call traces carry most of the signal. The verdict
records the success criterion, the capture channel, the planted and the
observed flag, the outcome, and which captures are present.
Listing~\ref{lst:verdict} shows the verdict of the privilege-escalation
record discussed in Section~\ref{sec:examples}. The dataset is committed with
the record and addressed by the hashes in \file{MANIFEST.sha256}, so that the
reference captures survive even if the images can no longer be rebuilt.

\begin{lstlisting}[float=t,caption={The \file{verdict.json} of a functional record.},label={lst:verdict}]
{
  "cve": "CVE-2023-22809",
  "success": true,
  "criterion": "flag",
  "channel": "stdout",
  "flag_planted":  "CVEX{9b8bb059600aa99e8a28eca00bb4ddef}",
  "flag_observed": "CVEX{9b8bb059600aa99e8a28eca00bb4ddef}",
  "benign": { "pcap": true, "syscalls": true },
  "attack": { "pcap": true, "syscalls": true }
}
\end{lstlisting}

\subsection{Reproducibility}

The format pursues reproducibility at four levels. Builds are meant to be
offline and pinned: base images are referenced by digest, and the vulnerable
software is vendored into the record, so that the application layer can be
built without network access. The runtime is pinned: a static subnet and
static addresses, explicit readiness conditions, and a wall-clock limit. The
flag is fixed and committed, so that reference captures are stable while
still proving the exploitation. Finally, the reference dataset is committed
and content-addressed.

Not all of these properties hold yet for every record. Images are not yet
pinned by digest, and records converted from \genie still download the
vulnerable software at build time rather than vendoring it. Packet timing is
not reproducible in any case: payloads and the order of exchanges are, and
timing is normalized away when datasets are compared.

\subsection{Limitations}

Because containers share the host kernel, the format cannot reproduce
vulnerabilities in the kernel itself, or in operating systems other than
Linux; Section~\ref{sec:future} discusses an extension to virtual machines
that shares the same manifest and dataset format. TLS interception assumes
that the client and target trust the recorder's certificate authority, which
is acceptable for research reproductions and is noted in each dataset.
Finally, the realism of the benign phase is defined by the record's author:
the value of the benign baseline depends on how representative the benign
workload is, and reviewers should check that it exercises the same interface
as the attack.

\section{Examples}
\label{sec:examples}

We describe two functional records from the corpus. Both concern well-known,
publicly disclosed and long-patched vulnerabilities; the point of the records
is not the vulnerabilities themselves, but the reusable, verified datasets
they produce. The first, a local privilege escalation in \texttt{sudo}, was
migrated by hand from the previous generation of \cvex and is verified by a
planted flag. The second, an HTTP request-smuggling flaw in the
\texttt{aiohttp} Python framework, was converted from a \genie reproduction
and is verified by \genie's CTF-style verifier. Table~\ref{tab:examples}
summarizes them.

\begin{table}[t]
\centering
\caption{The two example records.}
\label{tab:examples}
\small
\begin{tabular}{@{}lll@{}}
\toprule
 & \textbf{CVE-2023-22809} & \textbf{CVE-2024-52304} \\
\midrule
Software & \texttt{sudo} 1.9.12p1 & \texttt{aiohttp} 3.10.10 \\
CWE & CWE-269 & CWE-444 \\
Class & local privilege escalation & HTTP request smuggling \\
Origin & previous \cvex generation & \genie \\
Topology & \texttt{single} & \texttt{single} \\
Success criterion & planted flag & \genie verifier \\
Flag location & \texttt{/root/flag} (mode 0600) & --- \\
Traced process & \texttt{sudo} & \texttt{python} \\
Attack capture & 2,164 syscalls & 515~KB syscalls, 4.5~KB pcap \\
\bottomrule
\end{tabular}
\end{table}

\subsection{Local: sudoedit (CVE-2023-22809)}

\paragraph{The vulnerability.}
\texttt{sudoedit} (\texttt{sudo -e}) lets an administrator allow a user to
edit specific files with elevated privileges, without granting a root shell.
According to the public advisory, versions of \texttt{sudo} up to 1.9.12p1
mishandle the editor selected through the user's environment, so that a user
who is permitted to edit one file can edit other files as
well~\cite{cve}. Since the user controls which files are opened with
elevated privileges, the flaw amounts to a full privilege escalation. It was
fixed in \texttt{sudo} 1.9.12p2.

\paragraph{The target.}
The target image builds the vulnerable \texttt{sudo} 1.9.12p1 from a
vendored source tarball, so that the record does not depend on a
distribution package that may change or disappear. It creates an
unprivileged user and installs a sudoers policy that grants this user a
single right: to edit \file{/etc/motd} with \texttt{sudoedit}. At container
start, \file{plant-flag.sh} writes the record's flag to \file{/root/flag},
owned by root with mode 0600. The policy is the whole point of the record:
under a correct \texttt{sudo}, nothing the user is allowed to do can reveal
the flag.

\paragraph{Benign and attack phases.}
The benign phase exercises exactly the permission the policy grants: the
user edits \file{/etc/motd} through \texttt{sudoedit}, with an editor that
leaves the file unchanged. The attack phase starts from a fresh container
and runs, as the same unprivileged user, the publicly documented
exploitation of the flaw; it succeeds if it prints the flag. Both phases run
the same binary under the same policy, so the two traces differ only in what
the exploitation adds.

\paragraph{Verification and dataset.}
The record declares \texttt{success: flag} with the \texttt{stdout} channel:
the runner compares the flag printed by the attack with the planted one, and
in the reference run they match. Because the flag is readable only by root,
a match proves that the escalation genuinely happened, and not merely that
an exploit ran. The record is local, so its dataset is centered on system
calls: the recorder traces \texttt{sudo} and its children, producing 1,699
lines in the benign phase and 2,164 in the attack phase. The difference is
visible directly in the traces: the attack trace contains the access to
\file{/root/flag} and additional process executions, while the benign trace
never touches the flag. This pairing (same program, same policy, one benign
and one malicious use) is the kind of labeled data that host-based detectors
need.

\subsection{Network-Facing: aiohttp (CVE-2024-52304)}

\paragraph{The vulnerability.}
\texttt{aiohttp} is a widely used asynchronous HTTP client/server framework
for Python. According to its advisory, before version 3.10.11 the
pure-Python HTTP parser incorrectly handled newlines inside chunk extensions
of chunked request bodies~\cite{cve}. When the C extensions are not in use,
a front-end proxy and an \texttt{aiohttp} back end can disagree on where a
request ends, which enables HTTP request smuggling (CWE-444) and can let an
attacker bypass protections enforced by the proxy.

\paragraph{From a \genie reproduction to a record.}
\genie reproduced this CVE automatically, producing an exploit and a
verifier. The verifier forces the pure-Python parser, starts an
\texttt{aiohttp} server, runs the exploit against it, and checks whether the
server received two distinct requests although the client sent only one;
if so, it prints \genie's flag. The conversion to \cvexthree packages these
artifacts into a record: the target image is \texttt{python:3.11-slim} with
\texttt{aiohttp} 3.10.10 (the last release before the fix) installed from
PyPI, and the exploit and verifier are copied into it. The record's
\file{\_source} directory keeps \genie's original artifacts and the CVE
metadata (advisory, CWE, patch commit), so that the record's provenance is
preserved.

\paragraph{Verification and dataset.}
The record declares \texttt{success: verifier}: the attack phase runs
\genie's verifier inside the target, and the runner accepts the
reproduction if the verifier succeeds; in the reference run it prints the
flag \texttt{3xploit66full}. Because the verifier starts its own server and
talks to it over the loopback interface, the recorder, which shares the
target's network namespace, captures the smuggling exchange itself (a
4.5~KB packet capture) together with 515~KB of system-call trace of the
Python processes. The record is a network-facing vulnerability reproduced
within a single container: no separate attacker container is needed, yet
the dataset still contains the malicious HTTP traffic.

\paragraph{Limitations.}
This record also shows where the corpus still needs work. Its benign phase
only imports the library, a placeholder rather than a realistic workload;
for most records converted from \genie, the benign workloads have yet to be
written so that they exercise the same interface as the attack (here, normal
chunked requests to the same server). Records converted from \genie also
download their source at build time rather than vendoring it. Both are on
our list of improvements to the corpus.

\section{Future Work: Virtual Machines}
\label{sec:future}

Containers share the host's kernel. This is what makes \cvexthree records
light enough to run by the hundred on one machine, but it also bounds the
vulnerabilities the format can express. A record cannot install a different
kernel, load kernel modules, change kernel configuration, or run a non-Linux
operating system, and a privileged operation inside a container affects the
host as well. We plan to extend the format with a \emph{VM target}: a
manifest option that replaces the target container with a virtual machine
image, while keeping everything else in the record unchanged (the roles, the
success criteria, the two-phase run protocol, and the dataset layout). The
VM-based generation of \cvex (Section~\ref{sec:background}) already
demonstrated how to provision and instrument such targets; the extension
would bring that capability into the uniform \cvexthree record.

Several classes of vulnerabilities need this extension:

\begin{itemize}
\item \textbf{Kernel vulnerabilities.} Linux kernel flaws such as
CVE-2016-5195 (``Dirty COW''), CVE-2022-0847 (``Dirty Pipe''), and
CVE-2024-1086 (a use-after-free in \texttt{nf\_tables}) depend on a specific
vulnerable kernel version and configuration, which a container cannot
provide.

\item \textbf{Vulnerabilities in non-Linux systems.} Windows vulnerabilities
such as CVE-2017-0144 (the SMBv1 flaw known as EternalBlue), CVE-2019-0708
(``BlueKeep'', in the Remote Desktop service), and CVE-2021-34527
(``PrintNightmare'', in the Print Spooler) require a Windows guest.

\item \textbf{Container and isolation boundaries.} Vulnerabilities in the
container runtime itself, such as CVE-2019-5736 in \texttt{runc}, cannot be
reproduced safely inside the runtime they affect.

\item \textbf{Environment-dependent behavior.} Some user-space
vulnerabilities depend on how the C library and the kernel interact. While
building the corpus we were unable to reproduce CVE-2019-18276, a
\texttt{bash} privilege-management flaw, because it requires a C library and
kernel combination that preserves the saved user ID when \texttt{bash} drops
privileges, which our container environment does not provide.

\item \textbf{Records that need host privileges.} A few records in the
corpus already require additional container capabilities (for example,
permission to mount file systems or to configure traffic control) to
reproduce their vulnerability. Granting these weakens the isolation between
the record and the host; a VM target would restore it.
\end{itemize}

A VM target also enables richer observation. Besides packets and system
calls, a VM-based recorder can capture kernel events, Windows API calls, and
whole-system snapshots taken before and after the attack. The success
criteria of \cvexthree carry over unchanged: a flag readable only by root in
the guest, or a crash witness for a kernel panic, decides success just as it
does for containers.

\section{Related Work}
\label{sec:related}

\paragraph{Studies of reproducibility.}
Mu et al.~\cite{mu2018reproducibility} conducted the first large empirical
study of vulnerability reproducibility, measuring how often crowd-reported
memory-corruption vulnerabilities can be reproduced from public reports and
which missing information most often prevents it. Their findings motivate
formats, like \cvex, that capture the full environment of a reproduction
rather than a prose description of it.

\paragraph{Reproducible vulnerability datasets.}
ARVO~\cite{mei2026arvo} rebuilds thousands of memory-safety vulnerabilities
found by OSS-Fuzz in C/C++ projects, providing for each a build environment,
a triggering input, and the fixing patch. Magma~\cite{hazimeh2020magma}
forward-ports real bugs into recent library versions, with ground-truth
canaries, to benchmark fuzzers. Vulhub~\cite{vulhub} is a community
collection of Docker Compose environments for vulnerable applications. In
software engineering, Defects4J~\cite{just2014defects4j} established the
value of a curated, executable database of real faults for controlled
experiments. \cvex shares the goal of executable ground truth, but differs in
three respects: it is not limited to a single bug class or ecosystem, it
defines success uniformly through a planted flag, and it records paired
benign and attack activity at both the network and the system-call level.

\paragraph{Benchmarks for automated exploitation.}
The rise of LLM-based agents has produced several benchmarks built on real
vulnerabilities. CVE-Bench~\cite{zhu2025cvebench} packages 40 critical web
application CVEs in a sandbox that scores agents' exploitation attempts.
Cybench~\cite{zhang2025cybench} collects 40 professional Capture-the-Flag
tasks. CyberGym~\cite{wang2026cybergym} asks agents to generate PoCs for 1,507
vulnerabilities from 188 projects, given a description and the codebase, and
SEC-bench~\cite{lee2025secbench} automatically constructs PoC-generation and
patching tasks from real vulnerabilities. Earlier, the DARPA Cyber Grand
Challenge~\cite{cgc} and the AI Cyber Challenge~\cite{aixcc} evaluated
autonomous systems on purpose-built or injected vulnerabilities. These
benchmarks define environments and success checks for their own harnesses;
\cvex aims to be a harness-independent format, and its records can serve as
tasks for such benchmarks.

\paragraph{Automated reproduction.}
\genie~\cite{ullah2026cvegenie} automates the end-to-end reproduction of
CVEs, from collecting resources to building the project, exploiting it, and
verifying the result with a CTF-style verifier. \cvex is complementary: it
provides the portable, rerunnable packaging that turns an automated
reproduction into a durable artifact, and \NumGenieFunctional records of our
corpus were obtained by converting \genie's reproductions.

\paragraph{Vulnerability code datasets.}
A large body of work collects vulnerable and fixed code for training and
evaluating vulnerability detectors, including the Juliet test suite in
NIST's SARD~\cite{sard}, Devign~\cite{zhou2019devign}, BigVul~\cite{fan2020bigvul},
CVEfixes~\cite{bhandari2021cvefixes}, and PrimeVul~\cite{ding2025primevul}.
These datasets are static: they label code, but do not provide a way to run
the vulnerable program or to demonstrate exploitation. \cvex records can
complement them with executable evidence.

\paragraph{Intrusion detection datasets.}
Labeled traffic and system-call datasets are the traditional basis for
evaluating intrusion detection, from the DARPA 1999
evaluation~\cite{lippmann2000darpa} to ADFA-LD~\cite{creech2013adfa},
CICIDS2017~\cite{sharafaldin2018cicids}, and LID-DS~\cite{grimmer2023liddsr}.
Sommer and Paxson~\cite{sommer2010outside} discuss the pitfalls of
evaluating machine-learning detectors on such data. \cvex produces
per-vulnerability, paired benign/attack captures on demand, where the attack
half is guaranteed by the success criterion to contain a working
exploitation.

\paragraph{Kernel and OS-level reproduction.}
Reproducing kernel vulnerabilities requires control over the kernel itself.
syzbot~\cite{syzbot} continuously reproduces kernel bugs found by fuzzing in
virtual machines, and SyzBridge~\cite{zou2024syzbridge} adapts upstream PoCs
to downstream distribution kernels to assess their exploitability. These
systems motivate the VM extension of \cvex discussed in
Section~\ref{sec:future}.

\section{Conclusion}
\label{sec:conclusion}

A CVE entry records that a vulnerability exists; it does not let anyone
observe it. We presented \cvexthree, a standard format for executable CVEs
that packages the vulnerable target, a benign workload, an attack, a
recorder, and a machine-checkable success criterion into a self-contained,
container-based record. Inspired by Capture-the-Flag competitions and by the
verifiers of \genie, \cvexthree defines success as the capture of a secret
that only a working exploit can reach, and it runs every record twice, so
that each run yields a labeled pair of benign and attack captures at the
network and system-call level.

We release a corpus of \NumRecords records, \NumFunctional of which are
fully functional and ship with verified reference datasets. We hope that
\cvex will make reproductions durable and shareable artifacts, and that the
corpus will be useful to researchers who build detection models, assess
patches, and evaluate automated security tools. Functionalizing the
remaining records, and extending the format to virtual machines for
vulnerabilities that require OS-level control, are the next steps.


\bibliographystyle{ACM-Reference-Format}
\bibliography{refs}

\end{document}
